\subsection{TENSOR PRODUCT OF GRADED ALGEBRAS OF THE SAME TYPES}

Assuming the hypotheses of no. 7, Proposition 10 hold, suppose further that all the \(\Delta_{i}\) are equal to the same commutative monoid \(\Delta_{0}\) ; we can then consider on the tensor \(\varepsilon\) -product \({}^{\varepsilon}\bigotimes_{i\in I}E_{i}\) the total graduation of type \(\Delta_{0}\) , associated with the graduation of type \(\Delta = \Delta_{0}^{I}\) on this algebra (II, § 11, no. 1); we shall call \({}^{\varepsilon}\bigotimes_{i\in I}E_{i}\) , with this graduation, a graded tensor \(\varepsilon\) -product of type \(\Delta_{0}\) of the family \((E_{i})_{i\in I}\) of graded algebras of type \(\Delta_{0}\) .

Always preserving the notation of Proposition 10 of no. 7, suppose that F is also a graded A-algebra of type \(\Delta_0\) and that the \(f_i\) are homomorphisms of graded algebras of type \(\Delta_0\) . Then \(f: {}^{\varepsilon}\bigotimes_{i \in I} E_i \to F\) is also a homomorphism of graded algebras of type \(\Delta_0\) : for it follows from formula (30) (no. 7) that if \(x_i\) is homogeneous and of degree \(\alpha_i \in \Delta_0\) , \(\bigotimes_{i \in I} x_i\) and \(\prod_{i \in I} f_i(x_i)\) are both homogeneous of degree \(\sum_{i \in I} \alpha_i \in \Delta_0\) .

It can therefore be said that \(\bigotimes_{i\in I}E_i\) , with the total graduation of type \(\Delta_0\) , constitutes, together with the canonical mapping \(\psi\) , a solution of a third universal mapping problem, where \(\Sigma\) is the species of graded A-algebra of type \(\Delta_0\) , the morphisms are homomorphisms of graded algebras of type \(\Delta_0\) and the \(\alpha\) -mappings are the mappings \(\prod_{i}f_{i}\) , where, in addition to conditions (26) (of no. 7), it is assumed that each \(f_{i}\) is a homomorphism of graded algebras of type \(\Delta_0\) .

For every subset J of I, the canonical homomorphism \({}^{\varepsilon}\bigotimes_{i\in J}E_i\to {}^{\varepsilon}\bigotimes_{i\in I}E_i\) (no. 7) is, in fact, a homomorphism of graded algebras of type \(\Delta_0\) , as follows immediately from the above.

PROPOSITION 12. (``associativity'' of the tensor \(\varepsilon\) -product of graded algebras of the same types). With the notation of Proposition 10 of no. 7, suppose that all the \(\Delta_{i}\) are equal to the same monoid \(\Delta_{0}\) ; let \((J_{\lambda})_{\lambda \in L}\) be a partition of I. With the notation of Proposition 11 of no. 7, suppose that the right hand side of formula (32) (no. 7) depends only on the sums \(\alpha_{\lambda}^{\prime \prime} = \sum_{i \in J_{\lambda}} \alpha_{i}, \beta_{\mu}^{\prime \prime} = \sum_{j \in J_{\mu}} \beta_{j}\) , for every ordered pair \((\lambda, \mu)\) of distinct indices, all \(\alpha_{\lambda}^{\prime} \in \Delta_{\lambda}^{\prime}\) and all \(\beta_{\mu}^{\prime} \in \Delta_{\mu}^{\prime}\) ; let \(\varepsilon_{\lambda \mu}^{\prime \prime}(\alpha_{\lambda}^{\prime \prime}, \beta_{\mu}^{\prime \prime})\) denote the right hand side of (32). Then \((\varepsilon_{\lambda \mu}^{\prime \prime})\) is a system of commutation factors over the family \((\Delta_{\lambda}^{\prime \prime})_{\lambda \in L}\) , where \(\Delta_{\lambda}^{\prime \prime} = \Delta_{0}\) for all \(\lambda \in L\) . If \(E_{\lambda}^{\prime \prime}\) is the graded tensor \(\varepsilon\) -product of type \(\Delta_{0}\) of the family \((E_{i})_{i \in J_{\lambda}}\) , there exists one and only one isomorphism of graded algebras of type \(\Delta_{0}\) , \(w: {}^{\varepsilon}\bigotimes_{\lambda \in L} E_{\lambda}^{\prime \prime} \to {}^{\varepsilon}\bigotimes_{i \in I} E_{i}\) , such that

\[
w \left(\bigotimes_ {\lambda \in L} \left(\bigotimes_ {i \in J _ {\lambda}} x _ {i}\right)\right) = \bigotimes_ {i \in I} x _ {i}\tag{35}
\]

provided that total orderings are chosen on the \(\mathbf{J}_{\lambda}\) and on \(\mathbf{I}\) as described in no. 7, Proposition 11.

By the hypothesis, for \(\gamma\) , \(\delta\) in \(\Delta_0\) , \(\varepsilon_{\lambda \mu}''(\gamma, \delta) = \varepsilon_{i_0 j_0}(\gamma, \delta)\) for some \(i_0 \in J_\lambda\) and some \(j_0 \in J_\mu\) , as is seen by considering the elements \(\alpha_\lambda' = (\alpha_i)_{i \in J_\lambda}\) and \(\beta_\mu' = (\beta_j)_{j \in J_\mu}\) such that \(\alpha_{i_0} = \gamma\) , \(\alpha_i = 0\) for \(i \neq i_0\) , \(\beta_{j_0} = \delta\) , \(\beta_j = 0\) for \(j \neq j_0\) ; it follows immediately that the \(\varepsilon_{\lambda \mu}''\) form a system of commutation factors. The rest of the proof is then analogous to that of Proposition 11 (no. 7) and is left to the reader.

Note that the additional hypotheses of Proposition 12 are fulfilled when \(\Delta_0 = \mathbf{Z}\) and that \((\varepsilon_{ij})\) is, either the trivial system of factors \((\varepsilon_{ij}(\alpha_i,\beta_j) = 1\) for all \(i,j)\) , or the system of factors defined by \(\varepsilon_{ij}(\alpha_i,\beta_j) = (-1)^{\alpha_i\beta_j}\) ; in the latter case, the right hand side of formula (32) is equal to \((-1)^{\gamma}\) , where \(\gamma = \sum_{i\in J_{\lambda},j\in J_{\mu}}\alpha_{i}\beta_{j} = \left(\sum_{i\in J_{\lambda}}\alpha_{i}\right)\left(\sum_{j\in J_{\mu}}\beta_{j}\right)\) .

Remarks. (1) Let I be an infinite indexing set and \(\Delta_0\) a commutative monoid; let \((\Delta_i)_{i \in I}\) denote the family such that \(\Delta_i = \Delta_0\) for all \(i\) and suppose given for every ordered pair of distinct indices \((i,j)\) of I a mapping \(\varepsilon_{ij} : \Delta_i \times \Delta_j \to A\) satisfying conditions (22), (23) and (24) (no. 7); this will also be called a system of commutation factors over the family \((\Delta_i)\) . Consider a family \((E_i)_{i \in I}\) of graded A-algebras of type \(\Delta_0\) ; for each finite subset J of I, let \(E_J\) denote a graded tensor \(\varepsilon\) -product of type \(\Delta_0\) of the subfamily \((E_i)_{i \in J}\) (with an arbitrary choice of a total ordering on J). If J, \(J'\) are two finite subsets of I such that \(J \subset J'\) , a canonical homomorphism of graded algebras of type \(\Delta_0\) , \(h_{J^{\prime}J}:E_{J}\to E_{J^{\prime}}\) , has been defined above and the uniqueness properties of these homomorphisms show immediately that if \(J\subset J^{\prime}\subset J^{\prime\prime}\) are three finite subsets of I, then \(h_{J^{\prime\prime}J}=h_{J^{\prime\prime}J^{\prime}}\circ h_{J^{\prime}J}\) . Thus there is a direct system \((E_{J},h_{J^{\prime}J})\) of graded algebras of type \(\Delta_{0}\) (§3, no. 3), whose indexing set is the right directed set \(\mathfrak{F}(I)\) of finite subsets of I. The graded algebra of type \(\Delta_{0}\) , the direct limit of this direct system (§3, no. 3), is called a graded tensor \(\varepsilon\) -product of type \(\Delta_{0}\) of the family \((E_{i})_{i\in I}\) ; it is also denoted by \({}^{\varepsilon}\bigotimes_{i\in I}E_{i}\) . When all the \(\Delta_{i}\) are equal to \(\mathbf{Z}\) and \(\varepsilon_{ij}(\alpha_{i},\beta_{j})=(-1)^{\alpha_{i}\beta_{j}}\) , the tensor product \({}^{\varepsilon}\bigotimes_{i\in I}E_{i}\) is also called the skew

tensor product of the family \((\mathbf{E}_{i})_{i\in\mathbb{I}}\) and is denoted by \({}^{g}\bigotimes_{i\in I}E_{i}\) . We leave to the reader the task of formulating and proving the proposition which generalizes Proposition 10 of no. 7 to the case where I is infinite, as Proposition 8 of no. 5 generalizes Proposition 5 of no. 2 to the case where I is infinite. Note that the underlying A-module of \({}^{\varepsilon}\bigotimes_{i\in I}E_{i}\) is the same as that underlying the (non-graded) tensor product of the family \((\mathbf{E}_{i})_{i\in\mathbb{I}}\) of non-graded algebras defined in no. 5.

\begin{enumerate}
\def\labelenumi{(\arabic{enumi})}
\setcounter{enumi}{1}
\tightlist
\item
  Let E be a graded A-algebra of type \(\Delta_0\) (where \(\Delta_0\) is a commutative monoid) and \(\rho: \mathbf{A} \to \mathbf{B}\) a ring homomorphism; the graduation on \(\rho^*(\mathbf{E})\) (II, §11, no. 5) is identical with the graduation on the graded tensor product \(\mathbf{B} \otimes_{\mathbf{A}} \mathbf{E}\) , where B has the trivial graduation.
\end{enumerate}
% semantic-fragments: legacy-input-seam
\relax{}
